% !TeX root = ../../Book2.tex
\section{Marcinkiewicz 插值}
\label{b2:sec:1703}

弱型估计只控制输出超过某一高度的集合。若直接把它代入层蛋糕公式，端点处常会出现发散的高度积分。可以先把输入按大小分成两部分：大值部分交给较低的可积指数，小值部分交给较高的可积指数。分割高度随正在考察的输出高度改变，两个端点就分别控制同一个高度积分的两端。

\subsection{次线性算子}
\label{b2:subsec:170301}

设 \((X,\mathcal A,\mu)\)、\((Y,\mathcal B,\nu)\) 为测度空间，\(\mathbb K=\mathbb R\) 或 \(\mathbb C\)，\(\mathcal D\) 是 \(X\) 上有限值可测 \(\mathbb K\)-值函数的几乎处处等价类组成的向量空间。设 \(T\) 把 \(\mathcal D\) 映为 \(Y\) 上有限值可测函数的几乎处处等价类。若对所有 \(f,g\in\mathcal D\)、\(c\in\mathbb K\)，有
\[
 |T(f+g)|\leq |Tf|+|Tg|,
 \quad |T(cf)|=|c|\cdot|Tf|
 \quad\text{几乎处处},
\]
则称 \(T\) 为\term[cixianxing-suanzi]{次线性算子}[sublinear operator]。极大算子在 \(L^1+L^\infty\) 等已保证输出几乎处处有限的定义域上满足这些条件；线性算子以及线性算子的模也满足这些条件。这里保留的是输出大小的三角不等式，不要求输出本身能够线性叠加。

\ProseParagraphBreak

为什么把大值交给较小的指数？设 \(1\leq p_0<p<p_1<\infty\)，\(f\in L^p(X,\mu)\)，并取 \(a>0\)。令
\[
 f^{>a}=f\cdot\mathbf1_{\{|f|>a\}},
 \quad f^{\leq a}=f\cdot\mathbf1_{\{|f|\leq a\}}.
\]
在大值集合上，\(p_0-p<0\) 给出 \(|f|^{p_0-p}\leq a^{p_0-p}\)；小值集合上则有 \(|f|^{p_1-p}\leq a^{p_1-p}\)。所以
\[
 \int_{|f|>a}|f|^{p_0}\,\mathrm d\mu
 \leq a^{p_0-p}\|f\|_p^p,
 \quad
 \int_{|f|\leq a}|f|^{p_1}\,\mathrm d\mu
 \leq a^{p_1-p}\|f\|_p^p.
\]
若 \(p_1=\infty\)，后一条换成 \(\|f^{\leq a}\|_\infty\leq a\)。这也证明了每个中间 \(L^p\) 输入都属于和空间 \(L^{p_0}+L^{p_1}\)：两个截断确实落在算子的端点定义域中。

\ProseParagraphBreak

以上粗界适合检查截断的可积性，暂时不要把它们代入整个插值证明。这样代入后，两项都只剩下 \(t^{-p}\|f\|_p^p\) 一类的分布上界，层蛋糕积分又会遇到 \(\int\mathrm dt/t\)。要保留比弱型更多的信息，必须让集合 \(\{|f|>a\}\) 和 \(\{|f|\leq a\}\) 留在高度积分里面。

\subsection{弱型估计}
\label{b2:subsec:170302}

回顾\secref{b2:sec:1702}的记号，\(d_h(t)=\nu(\{|h|>t\})\)、\(t>0\)，表示目标函数的分布函数。若 \(T\) 在 \(L^{p_i}(X)\) 上满足弱型 \((p_i,p_i)\)，其含义是存在 \(A_i\geq0\)，使
\[
 d_{Th}(t)\leq
 \left(\dfrac{A_i\|h\|_{L^{p_i}(X)}}{t}\right)^{p_i}
 \quad(h\in L^{p_i}(X),\ t>0).
\]
当 \(p_i=\infty\) 时，使用的是强型条件
\(\|Th\|_{L^\infty(Y)}\leq A_i\|h\|_{L^\infty(X)}\)。

\ProseParagraphBreak

对定义域中的分解 \(f=f_0+f_1\)，如果 \(|Tf_0|\leq t/2\) 且 \(|Tf_1|\leq t/2\)，次线性便给出 \(|Tf|\leq t\)。因此，除一个零测集外，
\[
 \{|Tf|>t\}\subseteq
 \{|Tf_0|>t/2\}\cup\{|Tf_1|>t/2\}.
\]
由此可以分别使用两个端点，而不需要 \(Tf=Tf_0+Tf_1\)。若较高端点是 \(L^\infty\)，且其常数 \(A_1>0\)，取 \(a=t/(2A_1)\) 就使小值部分的输出不超过 \(t/2\)；此时只须估计大值部分的超水平集合。

\subsection{Marcinkiewicz 插值}
\label{b2:subsec:170303}

\term[marcinkiewicz-chazhi]{Marcinkiewicz 插值定理}[Marcinkiewicz interpolation theorem]先把两个同指数的弱型端点连接起来。这是后面从弱 \((1,1)\) 与强 \((2,2)\) 推出奇异积分中间指数估计时使用的版本。

\Needspace{6\baselineskip}
\begin{theorem}[Marcinkiewicz]\label{b2:thm:17-marcinkiewicz}
设 \((X,\mathcal A,\mu)\)、\((Y,\mathcal B,\nu)\) 为测度空间，\(\mathbb K=\mathbb R\) 或 \(\mathbb C\)，\(1\leq p_0<p_1\leq\infty\)。设 \(T\) 定义在 \(L^{p_0}(X;\mathbb K)+L^{p_1}(X;\mathbb K)\) 上，输出 \(Y\) 上可测函数的几乎处处等价类，且满足 \(|T(f+g)|\leq|Tf|+|Tg|\)、\(|T(cf)|=|c|\cdot|Tf|\) 几乎处处。设 \(A_0,A_1\geq0\)，并假设对每个有限端点 \(p_i\)，有
\[
 \nu(\{|Th|>t\})\leq
 \left(\dfrac{A_i\|h\|_{L^{p_i}(X)}}{t}\right)^{p_i}
 \quad(h\in L^{p_i}(X),\ t>0);
\]
若 \(p_1=\infty\)，该端点改为
\(\|Th\|_{L^\infty(Y)}\leq A_1\|h\|_{L^\infty(X)}\)。
对 \(p_0<p<p_1\)，以 \(1/\infty=0\) 约定唯一的 \(0<\theta<1\)，使
\[
 \frac1p=\frac{1-\theta}{p_0}+\frac\theta{p_1}.
\]
则对所有 \(f\in L^p(X;\mathbb K)\)，有
\[
 \|Tf\|_{L^p(Y)}
 \leq C_{p,p_0,p_1}A_0^{1-\theta}A_1^\theta\|f\|_{L^p(X)},
\]
其中常数只依赖所列指数。源和目标空间不要求无原子或 \(\sigma\)-有限。
\end{theorem}
\begin{proof}
先设 \(A_0,A_1>0\)，固定 \(f\in L^p(X)\)。这里 \(p<\infty\)，所以
\[
 E=\{|f|>0\}=\bigcup_{n\geq1}\{|f|>1/n\},
 \quad \mu(\{|f|>1/n\})\leq n^p\|f\|_p^p<\infty.
\]
因此 \(\mu\) 在 \(E\) 上的限制是 \(\sigma\)-有限的。以下与输入有关的非负重积分都支撑在 \(E\times(0,\infty)\) 上，Tonelli 定理可在那里使用；目标上的层蛋糕公式则用\cref{b2:thm:17-layer-cake}的一般测度版本。

先讨论 \(p_1<\infty\)。取一个待选择的 \(\kappa>0\)，在输出高度 \(t\) 处按输入高度 \(a=\kappa t\) 分割。前面已证明两个截断分别属于 \(L^{p_0}\)、\(L^{p_1}\)，故超水平集合估计给出
\begin{equation}\label{b2:eq:17-height-split}
 \begin{aligned}
 d_{Tf}(t)\leq{}&\left(\frac{2A_0}{t}\right)^{p_0}
       \int_{|f|>\kappa t}|f|^{p_0}\,\mathrm d\mu\\
 &+\left(\frac{2A_1}{t}\right)^{p_1}
       \int_{|f|\leq\kappa t}|f|^{p_1}\,\mathrm d\mu.
 \end{aligned}
\end{equation}
选择 \(a\) 与 \(t\) 成正比，使固定输入值 \(s=|f(x)|\) 对应的高度区间恰为 \(0<t<s/\kappa\) 或 \(t\geq s/\kappa\)。两个端点因而分别控制高度零点和无穷远处。

将\eqref{b2:eq:17-height-split}乘以 \(pt^{p-1}\)，对 \(t>0\) 积分。层蛋糕公式与 Tonelli 给出
\[
 \begin{aligned}
 \|Tf\|_p^p\leq{}&p(2A_0)^{p_0}
   \int_E |f|^{p_0}
      \left(\int_0^{|f|/\kappa}t^{p-p_0-1}\,\mathrm dt\right)\mathrm d\mu\\
 &+p(2A_1)^{p_1}
   \int_E |f|^{p_1}
      \left(\int_{|f|/\kappa}^\infty t^{p-p_1-1}\,\mathrm dt\right)\mathrm d\mu\\
 ={}&p\left(
   \frac{(2A_0)^{p_0}\kappa^{p_0-p}}{p-p_0}
  +\frac{(2A_1)^{p_1}\kappa^{p_1-p}}{p_1-p}\right)\|f\|_p^p.
 \end{aligned}
\]
第一段内积分有限恰好用到 \(p>p_0\)，第二段有限用到 \(p<p_1\)。左端在计算开始时允许为无穷，右端有限才证明 \(Tf\in L^p\)。若先取 \(\kappa=1\)，已经得到所求强型估计。

为了同时看清端点常数的作用，令两个分子中的常数因子相等，取
\[
 \kappa=\left(\frac{(2A_0)^{p_0}}{(2A_1)^{p_1}}\right)^{1/(p_1-p_0)}.
\]
由倒数指数关系，
\[
 p(1-\theta)=\frac{p_0(p_1-p)}{p_1-p_0},
 \quad p\theta=\frac{p_1(p-p_0)}{p_1-p_0}.
\]
于是两个分子均为
\(2^pA_0^{p(1-\theta)}A_1^{p\theta}\)，并得到显式的上界
\[
 \|Tf\|_p\leq
 2\left[p\left(\frac1{p-p_0}+\frac1{p_1-p}\right)\right]^{1/p}
 A_0^{1-\theta}A_1^\theta\|f\|_p.
\]
这项选择使两个分子相等，从而显示两个端点常数在上界中的幂次。

若 \(p_1=\infty\)，取 \(a=t/(2A_1)\)。小值部分的输出不超过 \(t/2\)，因此
\[
 d_{Tf}(t)\leq\left(\frac{2A_0}{t}\right)^{p_0}
       \int_{|f|>t/(2A_1)}|f|^{p_0}\,\mathrm d\mu.
\]
换序后，固定 \(x\) 的高度积分为
\(\int_0^{2A_1|f(x)|}t^{p-p_0-1}\,\mathrm dt\)。故
\[
 \|Tf\|_p^p\leq
 \frac{p(2A_0)^{p_0}(2A_1)^{p-p_0}}{p-p_0}\|f\|_p^p,
\]
即
\[
 \|Tf\|_p\leq
 2\left(\frac p{p-p_0}\right)^{1/p}
 A_0^{p_0/p}A_1^{1-p_0/p}\|f\|_p.
\]
此时 \(\theta=1-p_0/p\)，仍是同一个几何幂。最后，若某个 \(A_i=0\)，先用正的较大常数 \(A_i+\varepsilon\) 作上述估计，再令 \(\varepsilon\downarrow0\)，便得到相应结论。
\end{proof}

两个高度积分解释了为什么结论只在开区间内成立：当 \(p=p_0\) 或有限的 \(p=p_1\) 时，其中一个内积分成为对数发散。改变分割常数可以调整算子范数的上界，不能消除这一端点障碍。

\ProseParagraphBreak

定理假设算子已经在和空间上定义，因而截断可以直接送入算子。另一种常见情形是线性算子先只定义在有限测度支集的简单函数上，并在那里满足同样的两个端点界。这个类对高度截断封闭，证明先给出该类的中间强型界；它在有限指数 \(L^p\) 中稠密，且
\(\|Tf_n-Tf_m\|_p=\|T(f_n-f_m)\|_p\leq C\|f_n-f_m\|_p\)，
所以目标 \(L^p\) 的完备性产生唯一连续延拓。对一般次线性算子，若另外具有
\(|Tf-Tg|\leq|T(f-g)|\)，也可作同一构造。非负值次线性算子具有这项差值控制：分别把 \(f=g+(f-g)\)、\(g=f+(g-f)\) 代入次可加性即可。若只知道本节定义中的模不等式，输出的相位变化并不受控，不能由此断言算子本身连续或有唯一延拓。

\ProseParagraphBreak

输入和输出的指数也可以不同。下面保留一个适用于局部势核的版本；除了层蛋糕公式，还需要用到指数至少为一时的积分三角不等式。

\Needspace{6\baselineskip}
\begin{proposition}\label{b2:prop:17-off-diagonal}
设 \((X,\mathcal A,\mu)\)、\((Y,\mathcal B,\nu)\) 为测度空间，\(\mathbb K=\mathbb R\) 或 \(\mathbb C\)，\(1\leq r<\infty\)。设 \(T\) 定义在 \(L^1(X;\mathbb K)+L^\infty(X;\mathbb K)\) 上，输出 \(Y\) 上可测函数的几乎处处等价类，满足 \(|T(f+g)|\leq|Tf|+|Tg|\)、\(|T(cf)|=|c|\cdot|Tf|\) 几乎处处。假设存在 \(A\geq0\)、\(B>0\)，使
\[
 \begin{aligned}
 \nu(\{|Th|>t\})&\leq
 \left(\frac{A\|h\|_{L^1(X)}}t\right)^r
 &&(h\in L^1(X),\ t>0),\\
 \|Th\|_{L^\infty(Y)}&\leq B\|h\|_{L^\infty(X)}
 &&(h\in L^\infty(X)).
 \end{aligned}
\]
则对 \(1<p<\infty\)、\(f\in L^p(X;\mathbb K)\)，有
\[
 \|Tf\|_{L^{pr}(Y)}
 \leq 2\left(\frac p{p-1}\right)^{1/(pr)}
 A^{1/p}B^{1-1/p}\|f\|_{L^p(X)}.
\]
这里不要求两个测度空间 \(\sigma\)-有限。
\end{proposition}
\begin{proof}
固定 \(f\in L^p(X)\)，其非零支集 \(E\) 仍为 \(\sigma\)-有限。对输出高度 \(t>0\)，取
\[
 f_0=f\cdot\mathbf1_{\{|f|>t/(2B)\}},
 \quad f_1=f-f_0.
\]
有 \(\|f_0\|_1\leq(t/(2B))^{1-p}\|f\|_p^p\)、\(\|f_1\|_\infty\leq t/(2B)\)，所以两个端点都能使用，且小值输出不超过 \(t/2\)。于是
\[
 d_{Tf}(t)\leq\left(\frac{2A}t\right)^r
 \left(\int_E|f(x)|\cdot\mathbf1_{\{t<2B|f(x)|\}}\,\mathrm d\mu(x)\right)^r.
\]
置 \(q=pr\)，按层蛋糕公式，
\[
 \|Tf\|_q^q\leq q(2A)^r
 \int_0^\infty t^{q-r-1}
 \left(\int_E |f(x)|\cdot\mathbf1_{\{t<2B|f(x)|\}}\,\mathrm d\mu(x)\right)^r\mathrm dt.
\]
与同指数证明相比，输入积分现在还带有 \(r\) 次幂，不能直接把这次幂移进 \(x\) 积分。

记 \(\mathrm d\rho(t)=t^{q-r-1}\,\mathrm dt\)，并将括号内的非负函数记为 \(H(t)\)。在这个高度测度上使用积分 Minkowski 不等式，得到
\[
 \begin{aligned}
 \|H\|_{L^r(\rho)}
 &\leq\int_E |f(x)|
       \left(\int_0^{2B|f(x)|}t^{q-r-1}\,\mathrm dt\right)^{1/r}\mathrm d\mu(x)\\
 &=\frac{(2B)^{p-1}}{(q-r)^{1/r}}\|f\|_p^p.
 \end{aligned}
\]
这里 \(q-r=r(p-1)>0\) 保证零点处可积。所用的积分不等式也可以直接核验：\(r=1\) 时它就是 Tonelli；\(r>1\) 时，令 \(r'=r/(r-1)\)，对任意非负 \(h\in L^{r'}(\rho)\)、\(\|h\|_{r'}\leq1\)，先由 Tonelli 交换 \(x,t\) 积分，再对每个 \(x\) 用 Hölder，便有
\[
 \int_0^\infty H(t)\cdot h(t)\,\mathrm d\rho(t)
 \leq\int_E |f(x)|
       \left(\int_0^{2B|f(x)|}\mathrm d\rho(t)\right)^{1/r}\mathrm d\mu(x).
\]
为避免预设 \(H\in L^r(\rho)\)，先取
\(H_n=\min\{H,n\}\cdot\mathbf1_{[1/n,n]}\)。它有有限 \(L^r(\rho)\) 范数；非零时用
\(h=H_n^{r-1}/\|H_n\|_r^{r-1}\) 测试，得到同一个范数上界，再由单调收敛令 \(n\to\infty\)。所有重积分都在 \(E\times(0,\infty)\) 这个 \(\sigma\)-有限乘积空间上进行。

将 \(H\) 的估计代回，取 \(1/r\) 次方，得到
\[
 \|Tf\|_q^p
 \leq q^{1/r}(2A)\cdot\frac{(2B)^{p-1}}{(q-r)^{1/r}}\|f\|_p^p
 =2^p\left(\frac p{p-1}\right)^{1/r}AB^{p-1}\|f\|_p^p.
\]
再取 \(1/p\) 次方即得结论。若 \(A=0\)，可先用正的较大弱型常数，再取极限。
\end{proof}

这个指数可以在参数线上读出：对端点 \((1,r)\) 与 \((\infty,\infty)\)，取 \(\theta=1-1/p\)，则
\[
 \frac1q=\frac{1-\theta}r=\frac1{pr},
 \quad A^{1-\theta}B^\theta=A^{1/p}B^{1-1/p}.
\]
高度积分中的 Minkowski 论证使这个参数计算成为强型估计；假如 \(r<1\)，该积分三角不等式不再成立，以上证明便不能照用。

\ProseParagraphBreak

例如在 \(\mathbb R^d\) 上取 \(0<\beta<d\)，令
\[
 k(x)=|x|^{-\beta}\cdot\mathbf1_{\{0<|x|<1\}},
 \quad r=\frac d\beta>1,
 \quad \omega_d=|B(0,1)|.
\]
球坐标积分与超水平集合的半径给出
\[
 \|k\|_1=\frac{d\omega_d}{d-\beta},
 \quad
 d_k(s)=\begin{cases}\omega_d,&0<s<1,\\
                         \omega_ds^{-r},&s\geq1,
          \end{cases}
 \quad \|k\|_{r,\infty}=\omega_d^{1/r}.
\]
由\secref{b2:sec:1702}的集合积分范数
\(N_r(h)=\sup_{0<|F|<\infty}|F|^{-1/r'}\int_F|h|\)，有
\(N_r(k)\leq r'\omega_d^{1/r}\)。对 \(h\in L^1\)，把\cref{b2:prop:17-weak-complete}用于可积因子 \(h\) 与弱空间因子 \(k\)，得
\[
 \|k*h\|_{r,\infty}\leq N_r(k*h)
 \leq N_r(k)\cdot\|h\|_1.
\]
另一方面，\(k\in L^1\) 给出
\(\|k*h\|_\infty\leq\|k\|_1\cdot\|h\|_\infty\)。卷积积分在 \(L^1+L^\infty\) 上具有兼容的定义：可积部分的卷积几乎处处绝对收敛，有界部分的卷积处处绝对收敛；其和等于原输入的卷积积分几乎处处。于是命题给出
\[
 \|k*f\|_{pr}
 \leq 2\left(\frac p{p-1}\right)^{1/(pr)}
       (r'\omega_d^{1/r})^{1/p}
       \left(\frac{d\omega_d}{d-\beta}\right)^{1-1/p}\|f\|_p.
\]
取 \(d=3\)、\(\beta=1\)、\(p=2\)，就得到这个局部核的 \(L^2\to L^6\) 估计，两个端点常数可分别取 \(\tfrac32(4\pi/3)^{1/3}\) 与 \(2\pi\)。这是本命题给出的一种充分指数范围。若删除核的空间截断，输入 \(h=1\) 已使远场积分发散，所用的 \(L^\infty\) 端点便不成立；全空间分数积分需要另行估计远场。

\subsection{极大算子应用}
\label{b2:subsec:170304}

\Needspace{6\baselineskip}
\begin{corollary}\label{b2:thm:17-maximal-strong}
设 \(d\) 为正整数，\(1<p\leq\infty\)。在 \(\mathbb R^d\) 上取 Lebesgue 测度，\(M\)、\(M_u\) 分别为中心球和包含观察点的球上绝对平均的上确界；\(M_{\mathcal D}\) 为标准半开二进网格
\(\mathcal D=\{\,2^{-j}(k+[0,1)^d):j\in\mathbb Z,\ k\in\mathbb Z^d\,\}\)
上包含观察点的立方体绝对平均的上确界。则三个算子都把 \(L^p(\mathbb R^d;\mathbb C)\) 映入自身，并满足
\[
 \|Mf\|_p+\|M_uf\|_p+\|M_{\mathcal D}f\|_p
 \leq C_{d,p}\|f\|_p.
\]
各算子在 \(L^\infty\) 上的范数均不超过一。
\end{corollary}
\begin{proof}
\cref{b2:thm:17-maximal-weak}及 \(M_u\leq2^dM\) 给出前两个算子的弱 \((1,1)\) 界，\cref{b2:thm:17-dyadic-level}给出二进版本。每个绝对平均都不超过输入的本性上确界，故三个算子均为 \(L^\infty\) 压缩。在 \(L^1+L^\infty\) 上，可积部分的极大函数几乎处处有限，有界部分的极大函数处处有界；次可加性因而使三个算子都具有所需定义域。应用\cref{b2:thm:17-marcinkiewicz}的 \(p_0=1,p_1=\infty\) 情形，即得有限 \(p\) 的结论。
\end{proof}

对中心极大算子，还可保留一个以后用于端点问题的精细分布界。设 \(d\geq1\)、\(1\leq p<\infty\)、\(f\in L^p(\mathbb R^d)\)、\(t>0\)。分解
\(f=f\cdot\mathbf1_{\{|f|>t/2\}}+f\cdot\mathbf1_{\{|f|\leq t/2\}}\)，
大值部分可积，小值部分的极大平均不超过 \(t/2\)。对大值部分用弱型估计，得到
\begin{equation}\label{b2:eq:17-maximal-tail-refined}
 |\{Mf>t\}|\leq\frac{2\cdot5^d}{t}
                   \int_{|f|>t/2}|f|\,\mathrm dx.
\end{equation}
当 \(p>1\) 时，保留截断再作高度积分，便有
\[
 \|Mf\|_p^p\leq \frac{5^dp\,2^p}{p-1}\|f\|_p^p.
\]
分母 \(p-1\) 来自 \(\int_0^{2|f(x)|}t^{p-2}\,\mathrm dt\)。

\ProseParagraphBreak

这一常数在 \(p\downarrow1\) 时的增长，可以用具体输入检验。在实直线上取 \(f=\mathbf1_{[0,1]}\)。当 \(x>2\) 时，半径为 \(x\) 的中心区间给出 \(Mf(x)\geq1/(2x)\)，所以
\[
 \|Mf\|_p^p\geq2^{-p}\int_2^\infty x^{-p}\,\mathrm dx
 =\frac{2^{1-2p}}{p-1},
 \quad \|f\|_p=1.
\]
因此算子范数至少为一个固定常数乘以 \((p-1)^{-1/p}\)。在 \(1<p\leq2\) 上，这与 \((p-1)^{-1}\) 可比：若 \(\delta=p-1\in(0,1]\)，则
\[
 \frac{\delta^{-1/p}}{\delta^{-1}}
 =\delta^{\delta/(1+\delta)}\in[e^{-1/e},1].
\]
结合上界，实直线上的中心极大算子范数在 \(p\downarrow1\) 时确有 \(1/(p-1)\) 的量级。精确最佳常数还依赖所选的极大算子，不由这两个界确定。

\ProseParagraphBreak

极大函数还能同时控制一整族正则化算子，进而说明它们在 Lebesgue 点恢复原函数。下面的核可取复值，不要求非负。

\Needspace{6\baselineskip}
\begin{proposition}\label{b2:prop:17-approximate-identity}
设 \(d\) 为正整数，\(\psi\in L^1(\mathbb R^d;\mathbb C)\)、\(\int_{\mathbb R^d}\psi=1\)。假设存在 \(\varepsilon>0\)、\(0<C<\infty\)，使
\[
 |\psi(x)|\leq C(1+|x|)^{-d-\varepsilon}
 \quad\text{几乎处处}.
\]
对 \(t>0\)，置 \(\psi_t(x)=t^{-d}\cdot\psi(x/t)\)。若 \(1\leq p<\infty\)、\(f\in L^p(\mathbb R^d;\mathbb C)\)，则 \(\psi_t*f\to f\) 于 \(L^p\)。对每个满足 \(Mf(x)<\infty\) 的点，所有 \(t>0\) 的卷积积分均绝对收敛，且
\[
 \sup_{t>0}|\psi_t*f(x)|\leq C'\cdot Mf(x),
\]
其中 \(M\) 为中心 Hardy--Littlewood 极大算子，\(C'\) 只依赖 \(d,\varepsilon,C\)。对任一有限值可测代表 \(f\)，在它的每个 Lebesgue 点 \(x\) 都有 \(\psi_t*f(x)\to f(x)\) 当 \(t\downarrow0\)。
\end{proposition}
\begin{proof}
先说明共同的极大函数控制。把积分区域分为 \(|y|<t\) 及环带 \(2^{j-1}t\leq|y|<2^jt\)、\(j\geq1\)。内球上 \(|\psi_t|\leq Ct^{-d}\)，相应积分不超过 \(C\omega_d\cdot Mf(x)\)。第 \(j\) 个环带上，核的上界为
\(C2^{d+\varepsilon}t^{-d}2^{-j(d+\varepsilon)}\)，而 \(|f(x-y)|\) 的积分不超过
\(\omega_d(2^jt)^d\cdot Mf(x)\)。所以
\[
 \int_{\mathbb R^d}|\psi_t(y)|\cdot|f(x-y)|\,\mathrm dy
 \leq C_{d,\varepsilon}C
       \left(1+\sum_{j\geq1}2^{-j\varepsilon}\right)\cdot Mf(x).
\]
右端与 \(t\) 无关，且几何级数收敛。这给出所述上确界估计及全部尺度的绝对收敛；由弱型或强型极大估计，满足 \(Mf(x)<\infty\) 的点组成满测集。

范数收敛可以回用\cref{b2:thm:03-approximate-identity}：缩放保持 \(\int\psi_t=1\)、\(\|\psi_t\|_1=\|\psi\|_1\)，且对每个 \(a>0\)，
\(\int_{|y|\geq a}|\psi_t(y)|\,\mathrm dy
=\int_{|z|\geq a/t}|\psi(z)|\,\mathrm dz\to0\)。
其证明在本处也可直接写成
\[
 \|\psi_t*f-f\|_p
 \leq\int_{\mathbb R^d}|\psi(z)|\cdot
             \|f(\,\cdot-tz)-f\|_p\,\mathrm dz.
\]
有限指数的平移连续性使每个固定 \(z\) 的范数差趋零，而共同可积控制为
\(2\|f\|_p\cdot|\psi(z)|\)；支配收敛给出结论。

固定 \(f\) 的一个 Lebesgue 点 \(x\)，令
\[
 A_x(r)=\frac1{|B(x,r)|}\int_{B(x,r)}|f(y)-f(x)|\,\mathrm dy.
\]
有 \(A_x(r)\to0\) 当 \(r\downarrow0\)。此外 \(S_x=\sup_{r>0}A_x(r)<\infty\)：先选 \(r_0>0\) 控制 \(r<r_0\) 的平均；对 \(r\geq r_0\)，Hölder 给出
\[
 A_x(r)\leq (\omega_dr^d)^{-1/p}\|f\|_p+|f(x)|
 \leq(\omega_dr_0^d)^{-1/p}\|f\|_p+|f(x)|.
\]
这也说明 \(Mf(x)\leq S_x+|f(x)|<\infty\)，从而该点的卷积积分已有定义。由 \(\int\psi_t=1\)，将前面的环带估计用于 \(|f(x-y)-f(x)|\)，得到
\[
 |\psi_t*f(x)-f(x)|
 \leq C_{d,\varepsilon}C
       \left(A_x(t)+\sum_{j\geq1}2^{-j\varepsilon}A_x(2^jt)\right).
\]
给定误差要求，先利用 \(A_x\leq S_x\) 与几何级数选出一个固定的尾部起点，使尾和对所有 \(t\) 都小；剩下只有有限个半径，每个都随 \(t\downarrow0\) 趋零。因此右端趋零。
\end{proof}

这里恢复的是所选代表在 Lebesgue 点的取值。\cref{b2:thm:17-differentiation}保证这些点几乎处处存在；一个 \(L^p\) 等价类本身并不指定任意一点的函数值。
